\begin{itemize}[leftmargin=*]
  \item \textbf{RMSE did not improve.} No model in this report lowers the RTT RMSE of the previous ensemble by more than \MaxRMSEGain\% on any link (the largest drop is the Starlink DL ensemble). Section~\ref{sec:predictable} explains why: the level of the next RTT is close to unpredictable from 6\,s of input. Longer input windows, cross-link features (for example the state of the other paths), and explicit models of the Starlink 15\,s reconfiguration~\cite{tanveer2023making} are the next candidates.
  \item \textbf{Coverage under shift.} Conformal coverage holds exactly only for exchangeable data. Validation sessions are also used for early stopping, and sessions differ in time of day and load. The measured RTT coverage of \PICPMin--\PICPMax{} is close to the target but not guaranteed. Online adaptive conformal methods~\cite{gibbs2021adaptive} can correct the offset during operation.
  \item \textbf{iperf alignment.} The iperf3 reports are stamped at the start of each one-second interval (Section~\ref{sec:audit}). Re-stamping them at the end of the interval in the parsers is the clean fix.
  \item \textbf{One split, three seeds.} As in the initial study, a single session split is used. Differences of one or two percent (for example most RMSE changes) are within the seed spread and should not be read as real effects.
  \item \textbf{5G subsampling.} For the 5G links every fourth training window is used for all methods. The full set may help the 5G models slightly.
  \item \textbf{Scheduler evaluation.} The benefit for the multipath scheduler is measured in a trace-driven emulation (Section~\ref{sec:dup}). The satellite and 5G sessions were not recorded at the same time, so the emulation assumes independent paths. An end-to-end test of the final model inside the liquid tunnel is the main item for the next period.
  \item \textbf{Volatile OneWeb windows.} The final model widens its interval too little in the most volatile tenth of OneWeb windows and covers it at about 0.89 (Section~\ref{sec:why}). The anchored decoder also ranks the rare OneWeb DL spikes worse than the previous model. A calibration that depends on the regime, or a decoder that does not follow the last spike, are candidates.
  \item \textbf{Starlink UL shift.} Every interval under-covers the Starlink UL test sessions slightly (0.93 to 0.94), which points to a difference between the calibration and test sessions of that link that no method corrects.
  \item \textbf{Linear quantile model.} A linear quantile model with the same calibration comes close to the final model on interval score (Section~\ref{sec:why}). Where compute is tight, it is a reasonable deployment choice, at the cost of a worse point forecast.
\end{itemize}
